\chapter{Background}
\label{chap:background}

\noindent
This chapter sets out the technical basis the project depends on and the work it builds on. It moves from the security priorities of industrial control systems, through the use of software-defined networking to react to attacks in these environments, to the more specific problems of trusting the control plane, keeping security decisions stateful in the data plane, recognising the operation an attacker is attempting, and grading a response by the consequence of acting on an asset. The intent is that a reader who is technically competent but new to the topic can follow, and fairly assess, the design and evaluation presented in later chapters.

\section{Industrial control systems and their security priorities}
\label{sec:ics}

Industrial control systems (ICS) supervise and drive physical processes in sectors such as electricity, water, gas and manufacturing. A programmable logic controller (PLC) reads sensors, runs a control program and drives actuators such as pumps and valves, often within a cycle measured in milliseconds. This coupling to the physical world changes what security has to protect. The NIST guide to operational-technology security states that OT security objectives typically prioritise integrity and availability, followed by confidentiality, and must treat safety as an overarching priority \cite{nist80082}. Etxezarreta et al. give the same priority from the network side, noting in their survey that availability is a priority in industrial control systems, whose concerns include fault tolerance and human safety, so that an unexpected system outage is not acceptable \cite{etxezarreta2023survey}. A controller that is prevented from driving its process can therefore cause a physical hazard rather than only a disclosure of data. A defence for such an environment must therefore be judged not only on whether it stops an attack, but on whether the act of defending disturbs the process.

A brief note on how such a controller is programmed helps the reader who comes from networking rather than automation. A PLC runs on a repeating scan: each cycle it copies the field inputs into a process image, executes its logic, and writes the outputs. That logic lives in organisation blocks that the runtime calls automatically, such as OB30, the cyclic block that holds this project's tank control. Memory is addressed by area: \texttt{\%I} and \texttt{\%Q} are the input and output images at the bit level, \texttt{\%ID} and \texttt{\%QD} the 32-bit analogue channels, and \texttt{\%M} internal memory bits. So a level sensor might arrive at \texttt{\%ID100}, a valve be driven from \texttt{\%QD100}, and an operator button be held in a \texttt{\%M} bit. These conventions recur throughout the design and evaluation chapters.

\section{Software-defined networking and OpenFlow}
\label{sec:sdn}

Software-defined networking (SDN) separates the decision about how traffic is forwarded, the control plane, from the switches that carry it, the data plane. A central controller holds a global view of the network and installs forwarding rules into the switches, most commonly through the OpenFlow protocol, so the behaviour of the network can be changed in software rather than by reconfiguring each device by hand \cite{lara2014openflow}. For intrusion response this programmability is attractive, because when a threat is detected the controller can add or remove rules across the network in a coordinated way \cite{etxezarreta2023survey}. The same centralisation is also a weakness, discussed in Section~\ref{sec:controlplane}, since the controller becomes a single point through which the whole forwarding policy can be altered.

\section{Reactive intrusion response with SDN in ICS}
\label{sec:reactive}

Etxezarreta et al. survey how SDN has been used to build intrusion-response strategies for ICS and organise the field into a taxonomy of three families \cite{etxezarreta2023survey}. The first is dynamic traffic filtering, in which a detection technique such as deep packet inspection, signature matching, anomaly detection or machine learning identifies malicious traffic and the controller drops the offending packets or blocks the flow so that it does not reach the target device. The survey's own example of this family detects payload alteration on the path between a SCADA server and a PLC and drops the malicious control packets before they arrive \cite{etxezarreta2023survey}. The second family is network survivability, which uses the global view of the network to detect a failure or attack and reconfigure paths so that operation continues. The third is cyber deception, which includes moving-target defence, where addresses or paths are randomised to raise the attacker's uncertainty, and honeypots.

Two observations follow that motivate this project. First, the direct response actions in these strategies are to drop or to block, that is, to stop traffic reaching a device. Second, and more important, none of the three families conditions its response on the consequence of acting on the particular asset being protected. A blocked flow to a rogue scanner and a blocked flow to a PLC in the middle of a control loop are treated in the same way, even though the second can be as harmful as the attack it answers. This is the practical reason automated response is rarely enabled in operational technology: the obstacle is not the accuracy of detection but trust that the response will not itself trip the process. CARS addresses this gap by making the response conditional on the criticality of the protected asset and on the class of operation attempted, rather than applying a single uniform action.

\section{Trusting the control plane and its forwarding integrity}
\label{sec:controlplane}

Because the SDN controller sets the forwarding policy for the whole network, its compromise is a distinct and serious threat. Gardiner et al. show that the centralisation of control in SDN creates a single point of failure, and demonstrate controller-in-the-middle attacks in which a compromised or interposed controller manipulates the network in an ICS setting \cite{gardiner2021controller}. A response system that trusts its own control plane without question cannot claim to be safe against this class of attack.

The complementary defence is to verify the forwarding state itself. Melis et al. present a toolkit, built as plug-ins for the ONOS controller, that lets an administrator express security policies for an industrial SDN and check them formally against the installed rules. Their checker detects compromised forwarding caused by bogus injected flow-rules, inner loops and black-holes, which they note are hard to find through ordinary network scans, as well as the replacement or removal of legitimate rules \cite{melis2018policychecker}. The flow-integrity checker in CARS follows this line of work: it holds a trusted baseline of the security policy and reports any drift from it, so that tampering with the forwarding rules is caught even when the controller itself is the source of the change.

\section{Stateful security in the data plane}
\label{sec:stateful}

Early SDN security work observed that pushing every decision to the controller is itself exploitable. Shin et al. introduce AVANT-GUARD, which extends the OpenFlow data plane with two mechanisms. Connection migration has the switch complete the TCP handshake and involve the controller only once a connection is established, which sharply reduces the traffic sent to the control plane during scanning or denial-of-service attacks. Actuating triggers allow conditional actions to be taken within the data plane in response to network conditions \cite{shin2013avantguard}. The relevant principle for this project is that security decisions can carry connection state in the data plane rather than treating each packet in isolation. The stateful stage of the CARS pipeline follows that principle: it tracks connection state so that only properly established flows are admitted, and out-of-state packets that a stateless rule would pass are refused.

\section{Asset criticality and consequence-based prioritisation}
\label{sec:criticality}

The idea that a response should depend on the value of what is being protected is the basis of consequence-driven security engineering for control systems. The Idaho National Laboratory's Consequence-driven Cyber-informed Engineering examines high-consequence risk in critical infrastructure by first identifying the most damaging events an adversary could cause, then the key devices and components that enable them, the plausible cyber attack paths to those devices, and concrete mitigations and tripwires \cite{inl2016cce}. Operational guidance from CISA and international partners complements this at the inventory level. It recommends classifying operational-technology assets both by criticality, where a critical asset is one whose failure or compromise would have the most significant impact, and by function, such as control, monitoring, engineering or management \cite{cisa2025assetinventory}. CARS applies both ideas. Each protected asset is assigned a criticality tier according to the consequence of its compromise, and that tier, together with the class of operation, determines both the decision and the strength and duration of the response. The specific tiering used in this project is defined in Chapter~\ref{chap:execution}.

\section{Operation awareness and the false-data-injection threat}
\label{sec:opaware}

Traffic filtering in ICS can act on more than the source and destination of a flow. Because industrial protocols carry the operation being requested, deep packet inspection allows a defence to distinguish, for example, a read of a value from a write of one, and to treat an engineering download differently again; the survey lists deep packet inspection among the detection techniques used in the filtering family \cite{etxezarreta2023survey}. This matters because some of the most damaging ICS attacks are not floods but precise manipulations of control or sensor data. In a false-data-injection attack the adversary alters the values the controller reads or writes so that the process is driven to an unsafe state while its own readings appear normal; Etxezarreta et al. identify false-data injection as one of the main threats to ICS and propose a replica-based moving-target defence against it \cite{etxezarreta2026replica}. Reconnaissance is the usual first stage of such campaigns, and defending against it in ICS has itself been studied as a moving-target-defence problem \cite{samanis2025phd}. CARS recovers the operation from the traffic and uses it in the decision, so that a read from a PLC and a write to it are not treated alike.

\section{Industrial control system testbeds}
\label{sec:testbeds}

Evaluating ICS security work requires a realistic environment, and building one is itself a recognised research problem. Antonioli and Tippenhauer present MiniCPS, a toolkit built on the Mininet emulator for reproducible security research on cyber-physical system networks, supporting simulated devices and industrial protocols \cite{antonioli2015minicps}. Gardiner et al. observe that ICS security work often cannot be evaluated on real systems for safety and operational reasons, and distil design principles and lessons from building ICS security testbeds of their own \cite{gardiner2019oops}. This project takes the higher-fidelity option. Rather than a pure emulation, it uses a hardware testbed with two real Siemens PLCs driving a live physical process through the SDN fabric, so that a response can be judged against its actual effect on the process. The design of that testbed is described in Chapter~\ref{chap:execution}.
